\clearpage
\SourceLeaf{608}
\ArithmeticHead{579}{R}{INHALT.}
\begingroup
\newcommand{\ContentsRow}[5]{\par\noindent\hangindent=#1\hspace*{#1}#2\hspace{.4em}\mbox{\footnotesize #3}\hfill\mbox{#4\quad #5}\par}
\ContentsRow{1em}{IV.\quad Dritter Beweis des Fundamentalsatzes bei den quadratischen Resten in einer neuen Einkleidung}{.\,.\,.\,.\,.\,.\,.\,.\,.\,.\,.\,.\,.\,.\,.\,.\,.\,.\,.\,.\,.\,.\,.\,.\,.}{Seite}{33}
\ContentsRow{3.3em}{Bemerkung}{.\,.\,.\,.\,.\,.\,.\,.\,.\,.\,.\,.\,.\,.\,.\,.\,.\,.\,.\,.\,.\,.\,.\,.\,.}{—}{36}
\noindent\hspace*{.5em}\textls[120]{Kubische und biquadratische Reste.}\par
\ContentsRow{1em}{V.\quad Zwei Sätze über kubische und biquadratische Reste}{.\,.\,.\,.\,.\,.\,.\,.\,.\,.\,.\,.}{—}{37}
\ContentsRow{3.3em}{Bemerkung}{.\,.\,.\,.\,.\,.\,.\,.\,.\,.\,.\,.\,.\,.\,.\,.\,.\,.\,.\,.\,.\,.\,.\,.\,.}{—}{37}
\ContentsRow{1em}{VI.\quad Zur Lehre von den ganzen komplexen Zahlen}{.\,.\,.\,.\,.\,.\,.\,.\,.\,.\,.\,.\,.\,.\,.}{—}{38}
\ContentsRow{3.3em}{Bemerkung}{.\,.\,.\,.\,.\,.\,.\,.\,.\,.\,.\,.\,.\,.\,.\,.\,.\,.\,.\,.\,.\,.\,.\,.\,.}{—}{50}
\ContentsRow{1em}{VII.\quad Beweis eines Satzes aus der höhern Arithmetik zur Theorie der biquadratischen Reste gehörig}{.\,.\,.\,.\,.\,.\,.\,.\,.\,.\,.\,.\,.\,.\,.\,.\,.\,.\,.\,.\,.\,.\,.\,.\,.\,.\,.}{—}{51}
\ContentsRow{3.3em}{Bemerkung}{.\,.\,.\,\raisebox{.25em}{.}\,.\,.\,.\,.\,.\,.\,.\,.\,.\,.\,.\,.\,.\,.\,.\,.\,.\,.\,.\,.\,.}{—}{52}
\ContentsRow{1em}{VIII.\quad Zum Reziprozitätsgesetz der quadratischen und der biquadratischen Reste}{.\,.\,.\,.\,.}{—}{53}
\ContentsRow{3.3em}{Bemerkungen}{.\,.\,.\,.\,.\,.\,.\,.\,.\,.\,.\,.\,.\,.\,.\,.\,.\,.\,.\,.\,.\,.\,.\,.}{—}{54, 55}
\ContentsRow{1em}{IX.\quad Hauptmomente des Beweises für die biquadratischen Reste}{.\,.\,.\,.\,.\,.\,.\,.\,.\,.}{—}{56}
\ContentsRow{3.3em}{Bemerkung}{.\,.\,.\,.\,.\,.\,.\,.\,.\,.\,.\,.\,.\,.\,.\,.\,.\,.\,.\,.\,.\,.\,.\,.\,.}{—}{57}
\ContentsRow{1em}{X.\quad Zur Theorie der biquadratischen Reste}{.\,.\,.\,.\,.\,.\,.\,.\,.\,.\,.\,.\,.\,.\,.\,.\,.}{—}{58}
\ContentsRow{3.3em}{Bemerkung}{.\,.\,.\,.\,.\,.\,.\,.\,.\,.\,.\,.\,.\,.\,.\,.\,.\,.\,.\,.\,.\,.\,.\,.\,.}{—}{64}
\ContentsRow{1em}{XI.\quad Beweis des Reziprozitätssatzes für die biquadratischen Reste, der auf die Kreisteilung gegründet ist}{.\,.\,.\,.\,.\,.\,.\,.\,.\,.\,.\,.\,.\,.\,.\,.\,.\,.\,.\,.\,.\,.\,.\,.\,.}{—}{65}
\ContentsRow{3.3em}{Bemerkung}{.\,.\,.\,.\,.\,.\,.\,.\,.\,.\,.\,.\,.\,.\,.\,.\,.\,.\,.\,.\,.\,.\,.\,.\,.}{—}{69}
\noindent\hangindent=1.9em\hspace*{1em}XII.\quad Zur Geschichte der Theorie der kubischen und biquadratischen Reste.\par
\ContentsRow{3.3em}{\textsc{Gauss} an \textsc{Sophie Germain}, 1807 April 30}{.\,.\,.\,.\,.\,.\,.\,.\,.\,.\,.\,.\,.}{—}{70}
\ContentsRow{3.3em}{\textsc{Gauss} an \textsc{Olbers}, 1807 Juli 21}{.\,.\,.\,.\,.\,.\,.\,.\,.\,.\,.\,.\,.\,.\,.\,.\,.}{—}{74}
\ContentsRow{3.3em}{\textsc{Gauss} an \textsc{Olbers}, 1816 März 21}{.\,.\,.\,.\,.\,.\,.\,.\,.\,.\,.\,.\,.\,.\,.\,.\,.}{—}{75}
\ContentsRow{3.3em}{\textsc{Gauss} an \textsc{Bessel}, 1816 Dezember 23}{.\,.\,.\,.\,.\,.\,.\,.\,.\,.\,.\,.\,.\,.\,.}{—}{76}
\ContentsRow{3.3em}{Bemerkung}{.\,.\,.\,.\,.\,.\,.\,.\,.\,.\,.\,.\,.\,.\,.\,.\,.\,.\,.\,.\,.\,.\,.\,.\,.}{—}{76}
\par\medskip\noindent\textls[120]{Zur Theorie der Formen.}\par
\ContentsRow{1em}{I.\quad Über Polygonalzahlen}{.\,.\,.\,.\,.\,.\,.\,.\,.\,.\,.\,.\,.\,.\,.\,.\,.\,.\,.\,.\,.\,.\,.\,.}{—}{78}
\ContentsRow{3em}{Bemerkung}{.\,.\,.\,.\,.\,.\,.\,.\,.\,.\,.\,.\,.\,.\,.\,.\,.\,.\,.\,.\,.\,.\,.\,.\,.}{—}{79}
\ContentsRow{1em}{II.\quad Über die Anzahl der Zerlegungen einer Primzahl in drei Quadrate}{.\,.\,.\,.\,.\,.\,.\,.}{—}{80}
\ContentsRow{3em}{Bemerkungen}{.\,.\,.\,.\,.\,.\,.\,.\,.\,.\,.\,.\,.\,.\,.\,.\,.\,.\,.\,.\,.\,.\,.\,.}{—}{83}
\ContentsRow{1em}{III.\quad Über ternäre quadratische Formen}{.\,.\,.\,.\,.\,.\,.\,.\,.\,.\,.\,.\,.\,.\,.\,.\,.\,.\,.}{—}{86}
\ContentsRow{3em}{Bemerkungen}{.\,.\,.\,.\,.\,.\,.\,.\,.\,.\,.\,.\,.\,.\,.\,.\,.\,.\,.\,.\,.\,.\,.\,.}{—}{90}
\ContentsRow{1em}{IV.\quad Zur Bestimmung der Klassenanzahl}{.\,.\,.\,.\,.\,.\,.\,.\,.\,.\,.\,.\,.\,.\,.\,.\,.\,.\,.}{—}{91}
\ContentsRow{3em}{Bemerkung}{.\,.\,.\,.\,.\,.\,.\,.\,.\,.\,.\,.\,.\,.\,.\,.\,.\,.\,.\,.\,.\,.\,.\,.\,.}{—}{91}
\ContentsRow{1em}{V.\quad Zur Theorie der Formen}{.\,.\,.\,.\,.\,.\,.\,.\,.\,.\,.\,.\,.\,.\,.\,.\,.\,.\,.\,.\,.\,.\,.}{—}{92}
\ContentsRow{3em}{Bemerkung}{.\,.\,.\,.\,.\,.\,.\,.\,.\,.\,.\,.\,.\,.\,.\,.\,.\,.\,.\,.\,.\,.\,.\,.\,.}{—}{95}
\begin{center}\rule{.13\linewidth}{.4pt}\end{center}
\begin{center}\bfseries ALGEBRA.\end{center}
\noindent\emph{Briefwechsel.}\par
\noindent\hspace*{.5em}\textls[120]{Zum Fundamentalsatz der Algebra.}\par
\ContentsRow{2em}{1.\ \textsc{Pfaff} an \textsc{Gauss}, 1799 Mai 30}{.\,.\,.\,.\,.\,.\,.\,.\,.\,.\,.\,.\,.\,.\,.\,.\,.\,.\,.}{—}{99}
\Needspace{4\baselineskip}
\ContentsRow{2em}{2.\ \textsc{Pfaff} an \textsc{Gauss}, 1799 Juli 8}{.\,.\,.\,.\,.\,.\,.\,.\,.\,.\,.\,.\,.\,.\,.\,.\,.\,.\,.\,.}{—}{101}
\begin{flushright}73\textsuperscript{*}\end{flushright}
\endgroup
